% !TEX root = ../main.tex
\section{Enhancements}\label{sec:enhancements}

The base system described in \cref{sec:system-design} operates with greedy search, flat mutation instructions, and no cross-experiment memory.
This section describes four tracks of enhancements that improve representation, search, knowledge transfer, and evaluation fidelity.
\Cref{tab:enhancements} summarizes all enhancements with their track, mechanism, and benefit.

% ==================================================================
\subsection{Track 1: Representation}\label{sec:track-representation}

\subsubsection{Strategy Manifest}\label{sec:strategy-manifest}

The baseline mutation instruction is a monolithic \texttt{program.md} file.
The strategy manifest (\texttt{Strategy\allowbreak{}Manifest} in \texttt{anneal/\allowbreak{}engine/\allowbreak{}strategy.py}) replaces this with a structured YAML document containing four named components: \texttt{hypothesis\allowbreak{}\_generation}, \texttt{context\allowbreak{}\_assembly}, \texttt{mutation\allowbreak{}\_style}, and \texttt{failure\allowbreak{}\_analysis}.
Each component is a \texttt{Strategy\allowbreak{}Component} with an \texttt{approach} field (natural-language instructions), an \texttt{evolved\allowbreak{}\_at} counter tracking the last experiment at which it was rewritten, and a \texttt{streak\allowbreak{}\_without\allowbreak{}\_improvement} counter.

Component-level evolution targets the weakest link: \texttt{Strategy\allowbreak{}Manifest.\allowbreak{}weakest\allowbreak{}\_component()} returns the component with the longest streak without improvement.
The \texttt{evolve\allowbreak{}\_weakest\allowbreak{}\_component} function constructs a focused rewriting prompt that includes the current approach text, recent experiment statistics (kept/discarded counts, score ranges), and per-criterion pass/fail summaries.
The agent rewrites only the targeted component, leaving the remaining three stable.
This decomposition avoids catastrophic forgetting of working strategies when only one aspect of the mutation policy is underperforming.

The manifest also records a \texttt{lineage} (list of SHA hashes for experiments that triggered evolution) and an \texttt{aggregate\_score\_delta}, providing a traceable record of strategy drift.
Migration from \texttt{program.md} is automatic: \texttt{migrate\_program\_md} converts the monolithic content into the \texttt{hypothesis\_generation} component of a new manifest.

\subsubsection{Token-Efficient Context}\label{sec:context-compression}

Agent context grows with experiment history. The \texttt{context\_compression} field on \texttt{AgentConfig} supports three modes:
\begin{itemize}
    \item \texttt{none}: full context (experiment records, artifact content, knowledge).
    \item \texttt{moderate}: recent experiment records are truncated to hypothesis and outcome; full artifact content is retained.
    \item \texttt{aggressive}: only the last $N$ hypotheses and their outcomes are included; artifact content is summarized; duplicate information across knowledge context and experiment history is deduplicated.
\end{itemize}
Compression reduces input token cost and keeps the agent within its context window budget (\texttt{max\_context\_tokens}) without discarding the most recent signal.

% ==================================================================
\subsection{Track 2: Search}\label{sec:track-search}

\subsubsection{Dual-Agent Mutation via Thompson Sampling}\label{sec:dual-agent}

The \texttt{Agent\allowbreak{}Selector} class (\texttt{anneal/\allowbreak{}engine/\allowbreak{}strategy\allowbreak{}\_selector.py}) maintains two Beta-Binomial arms---\texttt{primary} and \texttt{exploration}---and selects between them using Thompson Sampling.
The primary agent uses the model and temperature specified in \texttt{Agent\allowbreak{}Config}; the exploration agent uses a separate model (\texttt{exploration\allowbreak{}\_model}) at higher temperature, biased toward novel mutations.

Selection is modulated by experiment progress: early in the run (\texttt{experiment\_ratio} $< 0.2$), exploration receives a $1.4\times$ boost; late in the run ($> 0.8$), exploration is dampened to $0.4\times$ its posterior sample.
After each experiment, the selected arm's posterior is updated with a binary reward (whether the mutation was kept).
This produces an adaptive explore-exploit balance without requiring manual scheduling.

\subsubsection{Two-Phase Mutation}\label{sec:two-phase}

When \texttt{two\_phase\_mutation} is enabled on \texttt{AgentConfig}, each experiment begins with a diagnosis step before the mutation agent is invoked.
The diagnosis call (\texttt{AgentInvoker.diagnose}) receives the current artifact, the most recent evaluation result (including per-criterion scores), and the last five experiment records.
It produces a \texttt{DiagnosisResult} (\texttt{anneal/engine/types.py}) containing:
\begin{itemize}
    \item \texttt{weakest\_criteria}: the criteria with the lowest scores.
    \item \texttt{root\_cause}: a natural-language explanation of why those criteria are failing.
    \item \texttt{fix\_category}: one of \texttt{structural}, \texttt{content}, \texttt{formatting}, \texttt{logic}, \texttt{coverage}, \texttt{other}.
    \item \texttt{suggested\_direction}: a concrete axis for the next mutation.
\end{itemize}
The diagnosis is prepended to the mutation prompt, focusing the agent on the identified weakness rather than generating undirected mutations.

\subsubsection{Simulated Annealing with Adaptive Reheat}\label{sec:simulated-annealing}

\texttt{SimulatedAnnealingSearch} (\texttt{anneal/engine/search.py}) accepts regressions with probability $\exp(\Delta / T)$ where $\Delta$ is the signed score delta and $T$ is the current temperature.
Temperature decays geometrically by a configurable \texttt{cooling\_rate} (default $0.95$) after each experiment, with a floor at \texttt{min\_temperature} (default $0.01$).

Adaptive reheat prevents premature convergence: when the acceptance ratio over the last \NUM{10} experiments drops below half the \texttt{acceptance\_target} (default $0.3$), the temperature is multiplied by \texttt{reheat\_factor} (default $2.0$), capped at the initial temperature.
This mechanism allows the search to escape local optima when the landscape becomes flat.

\texttt{HybridSearch} provides a sensible default for users who do not configure a search strategy: greedy for the first \NUM{10} experiments to establish a strong baseline, then simulated annealing to explore the neighborhood.

\subsubsection{Island Population Search}\label{sec:island-search}

\texttt{IslandPopulationSearch} (\texttt{anneal/engine/search.py}) maintains $N$ independent \texttt{PopulationSearch} instances, each with its own candidate pool and tournament selection.
Islands are assigned experiments in round-robin order.
Every \texttt{migration\_interval} experiments (default \NUM{10}), the best individual from each island is copied to every other island, injecting diversity without destroying each island's locally adapted population.

Each \texttt{PopulationSearch} instance supports tournament selection (sample $k$ candidates, discard the worst) and optional crossover: when \texttt{crossover\_rate} $> 0$, the two highest-scoring candidates' hypotheses are paired as parent inputs for the next mutation.

% ==================================================================
\subsection{Track 3: Knowledge}\label{sec:track-knowledge}

\subsubsection{Episodic Memory}\label{sec:episodic-memory}

After each experiment, \texttt{extract\_lesson} (\texttt{anneal/engine/knowledge.py}) produces a structured \texttt{Lesson} object (\texttt{anneal/engine/types.py}) containing:
\begin{itemize}
    \item \texttt{what\_changed}: a summary of the mutation diff.
    \item \texttt{what\_improved} / \texttt{what\_regressed}: per-criterion transitions (e.g., ``\texttt{clarity: PASS (was FAIL)}'').
    \item \texttt{transferable\_insight}: a domain-independent takeaway distilled by a lightweight LLM call.
    \item \texttt{domain\_tags}: tags for cross-domain retrieval (e.g., \texttt{error\_handling}, \texttt{code\_style}).
\end{itemize}

The \texttt{KnowledgeStore} persists experiment records in \texttt{experiments.jsonl} and maintains a hypothesis similarity index (TF-IDF with cosine similarity, upgraded to sentence embeddings when \texttt{sentence-transformers} is installed).
Context assembly follows a cold-start protocol: for the first \NUM{10} experiments, only raw records are injected; after that threshold, the agent receives the five most recent records, the five most similar past experiments by hypothesis embedding, and a consolidated learnings summary.

Every \NUM{50} experiments, deterministic consolidation produces a \texttt{ConsolidationRecord} summarizing kept/discarded/crashed counts, score trajectories, top improvements, failed approaches, tag frequencies, per-criterion variances, verifier block rates, and failure distribution.
Consolidation records are appended to \texttt{learnings-structured.jsonl} and rendered into a human-readable \texttt{learnings.md}.

\subsubsection{Lineage Context}\label{sec:lineage}

\texttt{KnowledgeStore.get\_lineage} traces the chain of \textsc{kept} mutations from the current best SHA back through \texttt{pre\_experiment\_sha} links, returning the last $d$ accepted ancestors (default $d = 5$).
This lineage is injected into the agent prompt as a causal chain: what was changed at each step and what score improvement resulted.
The lineage provides the agent with a narrative of how the current artifact reached its present state, avoiding mutations that revisit already-explored directions.

\subsubsection{Research Operator}\label{sec:research-operator}

When the optimization loop encounters a plateau (configurable consecutive experiments without improvement), the \texttt{ResearchOperator} (\texttt{anneal/engine/research.py}) injects external knowledge.
It uses a separate LLM call with a higher budget (\texttt{ResearchConfig.max\_budget\_usd}, default \$0.50) to generate up to \texttt{max\_suggestions} (default 3) novel mutation directions grounded in domain knowledge outside the experiment history.
Research hints are prepended to the next mutation prompt.
The operator auto-disables after \texttt{disable\_after\_failures} (default 3) consecutive failures to avoid wasting budget on unhelpful suggestions.

\subsubsection{Cross-Domain Knowledge Transfer}\label{sec:learning-pool}

The \texttt{LearningPool} (\texttt{anneal/engine/learning\_pool.py}) aggregates learnings extracted from experiments across targets and projects.
Each \texttt{Learning} is a frozen dataclass containing the observation text, signal polarity (\textsc{positive}/\textsc{negative}), score delta, per-criterion deltas, source condition, source target, and domain tags.

Retrieval is scope-aware (\texttt{LearningScope}: \texttt{condition}, \texttt{target}, \texttt{project}, \texttt{global}) and decay-adjusted: each learning's effective score is $|\Delta_{\text{score}}| \cdot e^{-\lambda t}$ where $\lambda$ is the decay rate (default $0.05$) and $t$ is the age in days.
Cross-domain learnings are penalized by a \texttt{domain\_penalty} factor (default $0.5$), prioritizing same-domain transfer.
The pool evicts lowest-impact learnings when it exceeds \texttt{max\_size} (default \NUM{1000}).

\texttt{PersistentLearningPool} backs the pool to \texttt{learnings-pool.jsonl} with file-level locking for concurrent safety.
\texttt{GlobalLearningPool} extends this to \texttt{\~{}/.anneal/global-learnings.jsonl}, enabling knowledge transfer across unrelated projects.

% ==================================================================
\subsection{Track 4: Evaluation}\label{sec:track-evaluation}

\subsubsection{Adaptive Sample Sizing}\label{sec:adaptive-sampling}

Fixed sample counts waste budget when the effect is large (more samples than needed) or miss real effects when the signal is weak (too few samples).
When \texttt{adaptive\_sampling} is enabled on \texttt{StochasticEval}, the evaluator starts with $\max(\texttt{min\_sample\_count},\; \lfloor\texttt{sample\_count} / 2\rfloor)$ samples and computes Cohen's $d$ against zero on the collected scores.
If $d > \texttt{early\_stop\_effect\_size}$ (default $1.0$), evaluation terminates early.
If $d < \texttt{extend\_effect\_size}$ (default $0.3$), up to $\lfloor\texttt{sample\_count} / 2\rfloor$ additional samples are collected, with a hard ceiling at $1.5 \times \texttt{sample\_count}$.
This mechanism reduces cost by \NUM{30}--\NUM{50}\% on clear-cut experiments while increasing statistical power on marginal ones.

\subsubsection{Eval Consistency Monitoring}\label{sec:drift-detection}

Evaluator drift---where the same artifact receives different scores over time---undermines the optimization signal.
The \texttt{KnowledgeStore.get\_drift\_report} method detects drift using a sliding-window comparison: within the most recent consolidation window, it splits records into first and second halves and flags any criterion whose variance increased by more than $2\times$ between halves, or whose total variance exceeds a threshold (default $0.1$).
Flagged criteria are reported as \texttt{DriftEntry} objects containing the criterion name, variance, mean score, and window size.

At the consolidation level, per-criterion variances are tracked in every \texttt{ConsolidationRecord}, producing a time series of evaluator stability.
Criteria exceeding the variance threshold trigger a warning log, alerting the operator to potential judge model degradation or prompt sensitivity.

\subsubsection{Eval Caching}\label{sec:eval-cache}

When the agent produces a mutation identical to a previously evaluated artifact, re-running the full evaluation pipeline wastes budget.
The \texttt{EvalCache} (\texttt{anneal/engine/eval\_cache.py}) is a content-addressed LRU cache keyed on the SHA-256 hash of the artifact content concatenated with the sorted criterion names.
On a cache hit, the stored \texttt{CacheEntry} (score, raw scores, criterion names) is returned immediately.

The cache also tracks \texttt{score\_history}---the sequence of scores observed for the same content hash across multiple evaluations.
The \texttt{consistency\_report} method flags entries where the standard deviation across evaluations exceeds $0.1$, surfacing potential evaluator inconsistency independent of the drift detection in the knowledge store.
Maximum cache size is configurable (default \NUM{200} entries) with LRU eviction.

% ==================================================================
% Summary table
% ==================================================================
\begin{table*}[t]
\centering
\caption{Summary of enhancements by track, mechanism, and benefit.}
\label{tab:enhancements}
\small
\begin{tabular}{@{}llllp{5.5cm}@{}}
\toprule
\textbf{ID} & \textbf{Name} & \textbf{Track} & \textbf{Mechanism} & \textbf{Benefit} \\
\midrule
E1 & Strategy Manifest & Representation & Structured YAML with named components & Component-level evolution; avoids catastrophic forgetting of working strategies \\
E2 & Token-Efficient Context & Representation & Three compression modes & Reduces input token cost; keeps agent within context budget \\
E3 & Dual-Agent Mutation & Search & Thompson Sampling over Beta-Binomial arms & Adaptive explore-exploit without manual scheduling \\
E4 & Two-Phase Mutation & Search & Diagnosis step before mutation & Focuses mutations on weakest criteria and root cause \\
E5 & Simulated Annealing & Search & $\exp(\Delta/T)$ acceptance with adaptive reheat & Escapes local optima; reheat prevents premature convergence \\
E6 & Island Population & Search & $N$ independent populations with migration & Maintains diversity; best individuals propagate across islands \\
E7 & Episodic Memory & Knowledge & Structured Lesson extraction + similarity retrieval & Avoids repeating failed approaches; transfers successful patterns \\
E8 & Lineage Context & Knowledge & Causal chain of accepted mutations & Agent understands how current artifact state was reached \\
E9 & Research Operator & Knowledge & External knowledge injection on plateau & Breaks out of exploration dead ends \\
E10 & Cross-Domain Transfer & Knowledge & Learning pool with decay and domain penalties & Transfers insights across targets and projects \\
E11 & Adaptive Sampling & Evaluation & Cohen's $d$ early-stop / extend & Reduces cost on clear experiments; increases power on marginal ones \\
E12 & Drift Detection & Evaluation & Sliding-window variance comparison & Detects evaluator degradation before it corrupts optimization \\
E13 & Eval Caching & Evaluation & Content-addressed LRU cache & Eliminates redundant evaluation of identical artifacts \\
\bottomrule
\end{tabular}
\end{table*}
